\section{Measuring the gain}
\label{sec:method}

\subsection{Probe items}

A \emph{probe item} freezes one decision point. It consists of a system prompt,
a goal, a prefix of steps that were executed successfully, a failing action
$\afail$, the observation the runtime returns for it, and the reference action
$\agold$ that would have succeeded. Scoring is teacher-forced: we never sample,
so the measurement has no decoding noise and is bit-reproducible.

Two properties of the item construction matter for the claims we make.

\paragraph{The failing action is fixed, not sampled.}
It would be natural to obtain $\afail$ by letting each model act and keeping its
mistakes. We deliberately do not: that yields a different item set per model, so
a scaling curve would confound model size with item difficulty, because larger models make
rarer, harder mistakes. Instead $\afail$ is produced from $\agold$ by a
fixed perturbation operator, so every model in the ladder is probed on
byte-identical items. The operators (Section~\ref{sec:setup}) are the mistakes
small models actually make: a dropped required argument, a pluralised tool name,
a renamed argument, a type confusion, a date written the way a person writes it,
a hallucinated entity.

\paragraph{Conditions differ in exactly one thing.}
For a given item, all conditions share the same system prompt, goal and
successful prefix, byte for byte. They differ only in how (or whether) the
failed attempt is written into the transcript. The conditions are:
\emph{pre} (no attempt recorded, the baseline of Eq.~\ref{eq:gain});
\emph{fail} (the attempt plus its error message, i.e.\ the standard harness);
\emph{succ} (the attempt plus a counterfactual observation reporting success);
\emph{neut} (the attempt plus a valence-free acknowledgement); and
\emph{abstract} (a harness-generated description of what went wrong, with the
call itself absent).

The counterfactual success observation is generated from the failing call's own
arguments, so it is coherent with the call it purports to answer. A generic
success string naming a different entity would introduce an incoherence cue that
the model could exploit, which would inflate $\Delta_{\mathrm{pol}}$ for reasons
having nothing to do with reading the error.

\subsection{Two behavioural readouts}

Log-probability differences are the primary measurement, but two derived
quantities are easier to reason about and we report both.

\paragraph{Normalised repeat probability.} Over the scored candidate set
$\mathcal{A} = \{\afail, \agold, d_1, d_2\}$, where $d_i$ are alternative wrong
actions for the same step,
$p_{\mathrm{rep}}(c) = \exp \log\pi(\afail \mid c) / \sum_{a \in \mathcal{A}} \exp \log \pi(a \mid c)$.
This is a proper probability over a fixed, model-independent action set and is
directly comparable across models.

\paragraph{Exact greedy repeat.} We record whether greedy decoding from a
context reproduces $\afail$ token for token. This is not an approximation: the
argmax at every position of the teacher-forced string is already computed while
scoring it, and if the argmax matches at all positions then greedy decoding
emits exactly that string. The statement ``the agent would have written the
failed call again, verbatim'' is therefore obtained at zero additional cost.

\subsection{Harness variants}
\label{sec:method:variants}

Section~\ref{sec:problem} argued that the remedy depends on which term dominates.
We evaluate four families of harness, all of which see the same environment
feedback and differ only in what reaches the model:

\begin{itemize}\itemsep2pt
\item \textbf{verbatim}: the standard ReAct transcript.
\item \textbf{+instruction}: verbatim, plus an explicit system-prompt
instruction not to repeat a call that has already failed.
\item \textbf{abstract}: the failed call is replaced by a description
generated by the runtime from its own error metadata, e.g.
      \texttt{[attempt 1 failed: create\_event, an argument was badly
formatted; that call is not repeatable]}. The diagnosis survives; the
      token sequence does not. No extra model call is involved.
\item \textbf{+ban}: a decoding-time constraint that makes
previously-failed action strings unreachable.
\end{itemize}

The ban is the classic bad-words construction: for each banned token sequence,
its final token is masked whenever the tokens generated so far match the
sequence's prefix (Algorithm~\ref{alg:ban}). Masking the \emph{first} token instead would be wrong, since it would forbid
every call beginning the same way,
including the corrected one, which typically shares a long prefix with the
failed call. We ban both the bare string and its leading-whitespace variant,
since byte-pair tokenisers assign these different first tokens.

\begin{algorithm}[t]
\small
\caption{Decoder-level suppression of failed actions}
\label{alg:ban}
\begin{algorithmic}[1]
\Require context $c$, ban list $B$ of token sequences, stop set $S$
\State $y \gets []$
\While{$|y| < L_{\max}$}
  \State $\ell \gets \log \pi(\cdot \mid c, y)$
  \For{$s \in B$}
    \State $k \gets |s| - 1$
    \If{$k = 0$ \textbf{or} $y_{-k:} = s_{1:k}$}
      \State $\ell[s_{|s|}] \gets -\infty$ \Comment{block only the completion}
    \EndIf
  \EndFor
  \State $y_{\text{next}} \gets \arg\max \ell$
  \State \textbf{if} $y_{\text{next}} \in S$ \textbf{then break}
  \State $y \gets y \Vert y_{\text{next}}$
\EndWhile
\State \Return $y$
\end{algorithmic}
\end{algorithm}

The ban costs one pass over $B$ per decoding step and no extra tokens, so in the
cost accounting of Section~\ref{sec:efficiency} it is free to three significant
figures. It is also \emph{exact}: a banned string cannot be emitted, so any
residual repetition must be a paraphrase. We measure that separately by
canonicalising actions (parsing the call and sorting its keyword arguments)
before counting repeats, which prevents the ban from looking better than it is.

\subsection{Efficient scoring}

The contexts we compare share a long prefix and differ in a short tail. Encoding
each independently would spend most of the compute re-encoding the system prompt
and the goal. We therefore keep one key/value cache together with the exact
token ids it holds; moving to a new context crops to the longest common prefix
and pushes only the remainder. On the ToolShed probe this removes
\CacheSaving\ of the token positions a cache-free implementation would
process. The same cache drives the agent rollouts, where step $t$'s context
extends step $t{-}1$'s, turning the per-step cost from ``re-encode the
transcript'' into ``encode the new turn''.

Because a silent numerical fault in this path would be indistinguishable from a
real effect, we re-score a random sample of items with a single cache-free
forward pass and require agreement to $10^{-3}$ nats; the check also rejects
non-finite and positive log-probabilities. Section~\ref{sec:limitations} explains
why that last guard is not paranoia.
